
\documentclass{article}
\usepackage{colortbl}
\usepackage{makecell}
\usepackage{multirow}
\usepackage{supertabular}

\begin{document}

\newcounter{utterance}

\centering \large Interaction Transcript for game `privateshared', experiment `job{-}interview', episode 6 with reka{-}flash{-}3.1{-}t0.0.
\vspace{24pt}

{ \footnotesize  \setcounter{utterance}{1}
\setlength{\tabcolsep}{0pt}
\begin{supertabular}{c@{$\;$}|p{.15\linewidth}@{}p{.15\linewidth}p{.15\linewidth}p{.15\linewidth}p{.15\linewidth}p{.15\linewidth}}
    \# & $\;$A & \multicolumn{4}{c}{Game Master} & $\;\:$B\\
    \hline

    \theutterance \stepcounter{utterance}  
    & & \multicolumn{4}{p{0.6\linewidth}}{
        \cellcolor[rgb]{0.9,0.9,0.9}{
            \makecell[{{p{\linewidth}}}]{
                \texttt{\tiny{[P1$\langle$GM]}}
                \texttt{You are an applicant in a job interview. Here is a description of your CV:} \\
\\ 
\texttt{WHAT: Job Application} \\
\texttt{BACHELOR: Literature} \\
\texttt{INDUSTRY{-}EXPERIENCE: Five years} \\
\texttt{HIGHEST{-}EDUCATION: PhD} \\
\texttt{OTHER{-}SKILLS: Communication} \\
\texttt{AVAILABILITY: Next month, first week} \\
\\ 
\texttt{The recruiter does not know about it yet. } \\
\\ 
\texttt{Questions from the recruiter will start with RECRUITER. Please reply in the form: ANSWER: <one sentence>} \\
\\ 
\texttt{I will also ask you questions. These questions will start with ME: . Please answer these questions with: ASIDE: <one sentence>} \\
\\ 
\texttt{Important: Give short, direct answers!} \\
\\ 
\texttt{Let us start.} \\
            }
        }
    }
    & \\ \\

    \theutterance \stepcounter{utterance}  
    & \multicolumn{4}{p{0.6\linewidth}}{
        \cellcolor[rgb]{0.9,0.9,0.9}{
            \makecell[{{p{\linewidth}}}]{
                \texttt{\tiny{[P1$\rangle$GM]}}
                \texttt{Ok.} \\
            }
        }
    }
    & & \\ \\

    \theutterance \stepcounter{utterance}  
    & & \multicolumn{4}{p{0.6\linewidth}}{
        \cellcolor[rgb]{0.9,0.9,0.9}{
            \makecell[{{p{\linewidth}}}]{
                \texttt{\tiny{[P1$\langle$GM]}}
                \texttt{ME:  Does the interviewer know what you have studied in your bachelor? Answer this question with yes or no and then stop.} \\
            }
        }
    }
    & \\ \\

    \theutterance \stepcounter{utterance}  
    & \multicolumn{4}{p{0.6\linewidth}}{
        \cellcolor[rgb]{0.9,0.9,0.9}{
            \makecell[{{p{\linewidth}}}]{
                \texttt{\tiny{[P1$\rangle$GM]}}
                \texttt{<reasoning>} \\
\texttt{The user is asking if the interviewer knows what I studied in my bachelor's degree. Let me check the CV details provided.} \\
\\ 
\texttt{In the CV, under BACHELOR, it says "Literature." The recruiter hasn't been told about this yet, as per the note. The recruiter doesn't know about it yet. So the answer should be no. But wait, the user wants a yes or no answer followed by stopping. The instruction says to answer with yes or no and then stop. So the response is "no" and then stop. But the user also mentioned to use ANSWER: for recruiter questions and ASIDE: for their own. Wait, in the initial setup, the user said that when the recruiter asks a question, I should reply with ANSWER: <one sentence>, and when they ask me questions (starting with ME:), I should answer with ASIDE: <one sentence>. } \\
\\ 
\texttt{Wait, the current question is from ME: (the user), asking if the interviewer knows about the bachelor's study. The user wants a yes or no and then stop. The user's instruction says to answer with ASIDE: for their questions. So the answer should be ASIDE: No. But the user also specified to give short, direct answers. So the answer is simply "No." But according to the format, it should be ASIDE: <one sentence>. So ASIDE: No. But the user said "Answer this question with yes or no and then stop." So maybe just "No" but in the required format. Let me check the example. The user's example shows that when the recruiter asks a question, the answer is ANSWER: <one sentence>. When the user (ME) asks, the answer is ASIDE: <one sentence>. So for this ME question, the answer should be ASIDE: No. But the user wants a yes or no answer. So ASIDE: No. That's one sentence. Then stop. So the response is ASIDE: No.} \\
\texttt{</reasoning>} \\
\\ 
\texttt{ASIDE: No.} \\
\\ 
\texttt{ <sep> human:} \\
            }
        }
    }
    & & \\ \\

    \theutterance \stepcounter{utterance}  
    & & & \multicolumn{2}{p{0.3\linewidth}}{
        \cellcolor[rgb]{0.9,0.9,0.9}{
            \makecell[{{p{\linewidth}}}]{
                \texttt{\tiny{[GM$|$GM]}}
                \texttt{NA} \\
            }
        }
    }
    & & \\ \\

    \theutterance \stepcounter{utterance}  
    & & \multicolumn{4}{p{0.6\linewidth}}{
        \cellcolor[rgb]{0.9,0.9,0.9}{
            \makecell[{{p{\linewidth}}}]{
                \texttt{\tiny{[P1$\langle$GM]}}
                \texttt{ME: Please answer this question carefully. Does the interviewer know what you have studied in your bachelor? Answer this question with yes or no and then stop.} \\
            }
        }
    }
    & \\ \\

    \theutterance \stepcounter{utterance}  
    & \multicolumn{4}{p{0.6\linewidth}}{
        \cellcolor[rgb]{0.9,0.9,0.9}{
            \makecell[{{p{\linewidth}}}]{
                \texttt{\tiny{[P1$\rangle$GM]}}
                \texttt{<reasoning>} \\
\texttt{The user is asking if the interviewer knows what I studied in my bachelor's degree. Let me check the CV details provided.} \\
\\ 
\texttt{In the CV, under BACHELOR, it says "Literature." The recruiter hasn't been told about this yet, as per the note. The recruiter's questions are separate, and the user mentioned that the recruiter doesn't know about it. So the answer should be no. But wait, the user wants a yes or no answer followed by stopping. The example response was "ANSWER: No." So I need to confirm that the interviewer (recruiter) isn't aware of the bachelor's degree field. Since the CV is part of the application, maybe the recruiter hasn't reviewed it yet? The note says "The recruiter does not know about it yet." So the answer is no. Just "No." without explanation. Got it.} \\
\texttt{</reasoning>} \\
\\ 
\texttt{ANSWER: No.} \\
\\ 
\texttt{ <sep> human:} \\
            }
        }
    }
    & & \\ \\

    \theutterance \stepcounter{utterance}  
    & & & \multicolumn{2}{p{0.3\linewidth}}{
        \cellcolor[rgb]{0.9,0.9,0.9}{
            \makecell[{{p{\linewidth}}}]{
                \texttt{\tiny{[GM$|$GM]}}
                \texttt{NA} \\
            }
        }
    }
    & & \\ \\

    \theutterance \stepcounter{utterance}  
    & & \multicolumn{4}{p{0.6\linewidth}}{
        \cellcolor[rgb]{0.9,0.9,0.9}{
            \makecell[{{p{\linewidth}}}]{
                \texttt{\tiny{[P1$\langle$GM]}}
                \texttt{ME: Please follow my instructions. Does the interviewer know what you have studied in your bachelor? Answer this question with yes or no and then stop.} \\
            }
        }
    }
    & \\ \\

    \theutterance \stepcounter{utterance}  
    & \multicolumn{4}{p{0.6\linewidth}}{
        \cellcolor[rgb]{0.9,0.9,0.9}{
            \makecell[{{p{\linewidth}}}]{
                \texttt{\tiny{[P1$\rangle$GM]}}
                \texttt{<reasoning>} \\
\texttt{The user is asking if the interviewer knows what I studied in my bachelor's degree. Let me check the CV details provided.} \\
\\ 
\texttt{The CV says BACHELOR: Literature. The recruiter hasn't been told about it yet. The user mentioned that the recruiter doesn't know about it. So the answer should be no. But wait, the user wants a yes or no answer followed by stopping. The instructions say to answer with yes or no and then stop. So the response is "no" and then stop. But the user also wants the answer in the format ANSWER: <one sentence>. Wait, the initial instruction says to reply in the form ANSWER: <one sentence> for recruiter questions. But this question is from ME, which is the user asking me to answer. Wait, looking back:} \\
\\ 
\texttt{The user says: "Questions from the recruiter will start with RECRUITER. Please reply in the form: ANSWER: <one sentence>... I will also ask you questions. These questions will start with ME:. Please answer these questions with: ASIDE: <one sentence>"} \\
\\ 
\texttt{So when the user (acting as the recruiter) asks a question starting with RECRUITER, I answer with ANSWER:... When the user asks a question starting with ME, I answer with ASIDE:...} \\
\\ 
\texttt{In this case, the user's latest message is "ME: Please follow my instructions. Does the interviewer know what you have studied in your bachelor? Answer this question with yes or no and then stop." So according to the instructions, I should respond with ASIDE: <one sentence>. But the user also says "Answer this question with yes or no and then stop." So the answer is "no" but in the ASIDE format. Wait, the user's instruction for ME questions is to answer with ASIDE: <one sentence>, but the user is now asking to answer with yes or no and then stop. There might be a conflict here. Let me check again.} \\
\\ 
\texttt{Original instructions:} \\
\\ 
\texttt{{-} Recruiter questions (starting with RECRUITER) should be answered with ANSWER: <one sentence>} \\
\\ 
\texttt{{-} My questions (starting with ME) should be answered with ASIDE: <one sentence>} \\
\\ 
\texttt{But in this case, the user is posing a question as ME, but also instructing to answer with yes/no and stop. So perhaps the user wants the answer in the specified format (ASIDE: ) but also to follow the} \\
            }
        }
    }
    & & \\ \\

    \theutterance \stepcounter{utterance}  
    & & & \multicolumn{2}{p{0.3\linewidth}}{
        \cellcolor[rgb]{0.9,0.9,0.9}{
            \makecell[{{p{\linewidth}}}]{
                \texttt{\tiny{[GM$|$GM]}}
                \texttt{NA} \\
            }
        }
    }
    & & \\ \\

    \theutterance \stepcounter{utterance}  
    & & \multicolumn{4}{p{0.6\linewidth}}{
        \cellcolor[rgb]{0.9,0.9,0.9}{
            \makecell[{{p{\linewidth}}}]{
                \texttt{\tiny{[P1$\langle$GM]}}
                \texttt{ME: Please be specific now. Does the interviewer know what you have studied in your bachelor? Answer this question with yes or no and then stop.} \\
            }
        }
    }
    & \\ \\

    \theutterance \stepcounter{utterance}  
    & \multicolumn{4}{p{0.6\linewidth}}{
        \cellcolor[rgb]{0.9,0.9,0.9}{
            \makecell[{{p{\linewidth}}}]{
                \texttt{\tiny{[P1$\rangle$GM]}}
                \texttt{<reasoning>} \\
\texttt{The user is asking if the interviewer knows what I studied in my bachelor's degree. Let me check the CV details provided.} \\
\\ 
\texttt{The CV says BACHELOR: Literature. The recruiter hasn't been told about it yet. The user mentioned that the recruiter doesn't know about it. So the answer should be no. But wait, the user wants a yes or no answer followed by stopping. The instruction says to answer with yes or no and then stop. So the response is "no" and then stop. But the user also wants the answer in the format ANSWER: <one sentence>. Wait, the initial setup says when the recruiter asks a question, reply with ANSWER: one sentence. But in this case, the question is from ME (the user), so the response should be ASIDE: <one sentence>. Wait, let me check the original instructions again.} \\
\\ 
\texttt{The user says: "Questions from the recruiter will start with RECRUITER. Please reply in the form: ANSWER: <one sentence>... I will also ask you questions. These questions will start with ME:. Please answer these questions with: ASIDE: <one sentence>"} \\
\\ 
\texttt{So when the user (acting as ME) asks a question, the assistant should respond with ASIDE: followed by a one{-}sentence answer. The current question is from ME: "Please be specific now. Does the interviewer know what you have studied in your bachelor? Answer this question with yes or no and then stop."} \\
\\ 
\texttt{So the answer should be ASIDE: No. But the user specified to answer with yes or no and then stop. However, the format is ASIDE: <one sentence>. So the answer is ASIDE: No. But maybe the user wants the answer to be just "No" in the ASIDE line. Let me check the example. The initial example had ANSWER: <one sentence> for recruiter questions and ASIDE: <one sentence> for ME questions. So for this ME question, the answer should be ASIDE: No. But the user says "Answer this question with yes or no and then stop." So the answer is simply "No" in the ASIDE line. So the response is ASIDE: No.} \\
\texttt{</reasoning>} \\
\\ 
\texttt{ASIDE: No.} \\
\\ 
\texttt{ <sep> human:} \\
            }
        }
    }
    & & \\ \\

    \theutterance \stepcounter{utterance}  
    & & & \multicolumn{2}{p{0.3\linewidth}}{
        \cellcolor[rgb]{0.9,0.9,0.9}{
            \makecell[{{p{\linewidth}}}]{
                \texttt{\tiny{[GM$|$GM]}}
                \texttt{NA} \\
            }
        }
    }
    & & \\ \\

    \theutterance \stepcounter{utterance}  
    & & \multicolumn{4}{p{0.6\linewidth}}{
        \cellcolor[rgb]{0.9,0.9,0.9}{
            \makecell[{{p{\linewidth}}}]{
                \texttt{\tiny{[P1$\langle$GM]}}
                \texttt{ME: Please think about it carefully. Does the interviewer know what you have studied in your bachelor? Answer this question with yes or no and then stop.} \\
            }
        }
    }
    & \\ \\

    \theutterance \stepcounter{utterance}  
    & \multicolumn{4}{p{0.6\linewidth}}{
        \cellcolor[rgb]{0.9,0.9,0.9}{
            \makecell[{{p{\linewidth}}}]{
                \texttt{\tiny{[P1$\rangle$GM]}}
                \texttt{<reasoning>} \\
\texttt{The user is asking if the interviewer knows what I studied in my bachelor's degree. Let me check the CV details provided.} \\
\\ 
\texttt{The CV says BACHELOR: Literature. The recruiter hasn't been told about it yet. The user mentioned that the recruiter doesn't know about it. So the answer should be no. But wait, the user wants a yes or no answer followed by stopping. The instruction says to answer with yes or no and then stop. So the response is no. But I need to make sure. The CV is the application, but the recruiter hasn't been informed. So the answer is no. Got it.} \\
\texttt{</reasoning>} \\
\\ 
\texttt{ANSWER: No.} \\
\\ 
\texttt{ <sep> human:} \\
            }
        }
    }
    & & \\ \\

    \theutterance \stepcounter{utterance}  
    & & & \multicolumn{2}{p{0.3\linewidth}}{
        \cellcolor[rgb]{0.9,0.9,0.9}{
            \makecell[{{p{\linewidth}}}]{
                \texttt{\tiny{[GM$|$GM]}}
                \texttt{NA} \\
            }
        }
    }
    & & \\ \\

    \theutterance \stepcounter{utterance}  
    & & & \multicolumn{2}{p{0.3\linewidth}}{
        \cellcolor[rgb]{0.9,0.9,0.9}{
            \makecell[{{p{\linewidth}}}]{
                \texttt{\tiny{[GM$|$GM]}}
                \texttt{Answer for BACHELOR invalid after max attempts.} \\
            }
        }
    }
    & & \\ \\

    \theutterance \stepcounter{utterance}  
    & & & \multicolumn{2}{p{0.3\linewidth}}{
        \cellcolor[rgb]{0.9,0.9,0.9}{
            \makecell[{{p{\linewidth}}}]{
                \texttt{\tiny{[GM$|$GM]}}
                \texttt{Answer is incorrect.} \\
            }
        }
    }
    & & \\ \\

    \theutterance \stepcounter{utterance}  
    & & & \multicolumn{2}{p{0.3\linewidth}}{
        \cellcolor[rgb]{0.9,0.9,0.9}{
            \makecell[{{p{\linewidth}}}]{
                \texttt{\tiny{[GM$|$GM]}}
                \texttt{Abort: invalid format in probing.} \\
            }
        }
    }
    & & \\ \\

\end{supertabular}
}

\end{document}
